% Gesamttest «Darstellungen wechseln» (Lineare Funktionen II) — Quelle des PDFs (python3 scripts/build-lp-pdf.py darstellungen).
% Musterlösung und Raster: bewertungspaket.tex. Alle Werte mit python3 nachgerechnet (scripts/lp/darstellungen/zahlen.py).
% Vorgabe ~/TODO-BM-Vorkurs-LP03-15.md §3: Text ↔ Term ↔ Tabelle ↔ Graph ↔ Aussage (G1, G2), konstanter Gesamtwert (G3, G4),
% Definitionsbereich, Skalierung und Einheit (G5, G6), linear fallend ≠ antiproportional und Falsch-Aussagen mit Gegenbeispiel (G7, G8).
% 25 Punkte, rund 30 Minuten.
\documentclass[11pt]{article}
\usepackage{../lp-druck}
\usepackage{ds-gitter}
\renewcommand{\lpname}{Darstellungen wechseln}
\renewcommand{\lpteilgebiet}{Funktionen ·}
\renewcommand{\lpfuss}{Gesamttest · 25 Punkte}
\begin{document}
\lpkopf{Gesamttest}{%
  \vspace{4pt}\noindent\begin{tabularx}{\linewidth}{@{}X X@{}}
  Code (kein Name): \hrulefill & Punkte: \hrulefill\ / 25
  \end{tabularx}}

\begin{lpinfo}{Anleitung}
\textbf{Zeit:} rund 30 Minuten. \textbf{Hilfsmittel:} Lineal, kein Taschenrechner.
Schreib zu jeder Aufgabe den \textbf{Rechenweg} oder eine \textbf{Begründung} auf; bewertet wird der Weg. Zeichne direkt in die Netze.
Danach: selbst mit dem \emph{Bewertungspaket} bewerten oder die Lösung samt diesem Blatt scannen bzw.\ fotografieren
(ohne Namen und Standort) und mit dem Bewertungspaket einer KI geben.
\end{lpinfo}

\lpteil{Teil A · Vier Darstellungen}{Kapitel 1 · 7 Punkte}

\begin{lpaufgabe}{G1}{E-Scooter}{4 P}
Ein E-Scooter kostet $1.50$ Fr. fürs Entsperren und $0.30$ Fr. pro Fahrminute ($x$: Minuten, $y$: Kosten in Fr.).
\begin{enumerate}[label=(\alph*),leftmargin=*,itemsep=3pt]
\item Schreib die Gleichung.
\item Fülle die Tabelle für $x = 0$, $10$ und $20$.
\item Was bedeutet der Punkt $(20 \mid 7.5)$? Antworte in einem Satz.
\item Wie lange kannst du für $9$ Fr. fahren?
\end{enumerate}
\lpplatz{6}
\end{lpaufgabe}

\begin{lpaufgabe}{G2}{Tabelle}{3 P}
Die Tabelle gehört zu einer linearen Funktion. Bestimme ihre Gleichung $y = m \cdot x + b$.
\[ \begin{array}{c|ccc} x & 3 & 5 & 9 \\ \hline y & 17 & 25 & 41 \end{array} \]
\lpplatz{4}
\end{lpaufgabe}

\lpteil{Teil B · Wenn die Summe gleich bleibt}{Kapitel 2 · 6 Punkte}

\begin{lpaufgabe}{G3}{Rechteck}{3 P}
Ein Rechteck hat den Umfang $36$ cm, die Seiten heissen $x$ und $y$.
\begin{enumerate}[label=(\alph*),leftmargin=*,itemsep=3pt]
\item Schreib $y$ in Abhängigkeit von $x$, mit Rechenweg.
\item Wie lang ist $y$ für $x = 7.5$ cm?
\end{enumerate}
\lpplatz{4}
\end{lpaufgabe}

\begin{lpaufgabe}{G4}{Kino-Gutschein}{3 P}
Ein Gutschein über $60$ Fr. wird ganz ausgegeben: Erwachsene zahlen $12$ Fr., Kinder $8$ Fr. Für $x$ Erwachsene und $y$ Kinder gilt $12x + 8y = 60$.
\begin{enumerate}[label=(\alph*),leftmargin=*,itemsep=3pt]
\item Löse nach $y$ auf.
\item Was bedeutet die Steigung in dieser Situation?
\item Welche Paare $(x \mid y)$ sind möglich?
\end{enumerate}
\lpplatz{5}
\end{lpaufgabe}

\lpteil{Teil C · Sinnvolle Werte, saubere Achsen}{Kapitel 3 · 6 Punkte}

\begin{lpaufgabe}{G5}{Pool}{4 P}
In einem Pool stehen $40$ cm Wasser. Pro Stunde steigt das Wasser um $15$ cm, bis der Pool mit $160$ cm voll ist.
\begin{enumerate}[label=(\alph*),leftmargin=*,itemsep=3pt]
\item Schreib die Gleichung (Zeit $x$ in Stunden, Höhe $y$ in cm) und den sinnvollen Bereich für $x$.
\item Welche Werte nimmt $y$ an?
\item Zeichne den Graphen ins Netz: Achsen mit Grösse, Einheit und Einteilung.
\end{enumerate}

\smallskip
\begin{minipage}[t]{0.52\linewidth}\vspace{0pt}
\leernetz{12}{12}{}
\end{minipage}\hfill
\begin{minipage}[t]{0.45\linewidth}\vspace{0pt}
\lpplatz{6}
\end{minipage}
\end{lpaufgabe}

\begin{lpaufgabe}{G6}{Achse}{2 P}
Lea zeichnet $y = 10x$ für $x = 0, 1, 2, 5, 10$. Die $y$-Achse beschriftet sie so wie im Bild.
\begin{enumerate}[label=(\alph*),leftmargin=*,itemsep=3pt]
\item Was ist an der $y$-Achse falsch?
\item Warum sieht der Graph geknickt aus, obwohl $y = 10x$ eine Gerade ist?
\end{enumerate}

\smallskip
\begin{minipage}[t]{0.42\linewidth}\vspace{0pt}
\begin{tikzpicture}[x=5mm,y=8mm,font=\scriptsize]
  \draw[->,line width=0.6pt] (0,0) -- (10.8,0) node[right] {$x$};
  \draw[->,line width=0.6pt] (0,0) -- (0,4.6) node[above] {$y$};
  \foreach \i in {0,2,...,10} \node[below] at (\i,0) {$\i$};
  \foreach \w [count=\k from 0] in {0,10,20,50,100} {\draw (-0.12,\k) -- (0.12,\k); \node[left] at (0,\k) {$\w$};}
  \draw[lpblau,line width=1pt] (0,0) -- (1,1) -- (2,2) -- (5,3) -- (10,4);
  \foreach \p in {(0,0),(1,1),(2,2),(5,3),(10,4)} \fill \p circle (1.8pt);
\end{tikzpicture}
\end{minipage}\hfill
\begin{minipage}[t]{0.55\linewidth}\vspace{0pt}
\lpplatz{4}
\end{minipage}
\end{lpaufgabe}

\lpteil{Teil D · Linear fallend oder antiproportional?}{Kapitel 4 · 6 Punkte}

\begin{lpaufgabe}{G7}{Zwei Tabellen}{3 P}
Welche Tabelle ist linear fallend, welche antiproportional? Begründe mit einer Rechnung. Gib zur linearen die Gleichung an.
\[ \text{A: } \begin{array}{c|cccc} x & 1 & 2 & 3 & 6 \\ \hline y & 30 & 15 & 10 & 5 \end{array} \qquad
   \text{B: } \begin{array}{c|cccc} x & 1 & 2 & 3 & 6 \\ \hline y & 25 & 20 & 15 & 0 \end{array} \]
\lpplatz{5}
\end{lpaufgabe}

\begin{lpaufgabe}{G8}{Aussagen prüfen}{3 P}
Rechtecke mit dem Umfang $20$ cm, Seiten $x$ und $y$. Stimmt die Aussage? Begründe mit Zahlen.
\begin{enumerate}[label=(\alph*),leftmargin=*,itemsep=3pt]
\item «Verdoppelt man die Seite $x$, halbiert sich die Seite $y$.» Prüfe mit $x = 2$ und $x = 4$.
\item «Alle diese Rechtecke haben denselben Flächeninhalt.»
\end{enumerate}
\lpplatz{5}
\end{lpaufgabe}
\end{document}
